\section{Open questions}\label{sec:questions}

\begin{enumerate}
\item Is $\Aut(X_{F_{54}})$ generated by $\sh$, $F$, the wall relabellings of Theorem~\ref{thm:54} and the wall-trading involutions of the kind in Theorem~\ref{thm:tau}, of all radii? At radius $2$ the relabellings suffice (Theorem~\ref{thm:54census}); at radius $3$ they do not, by Theorem~\ref{thm:tau}. A complete count of $\Cent^{\times}_3(F_{54})$ would be a first test.
\item Determine $\Aut(X_{F_{90}})$. Its elements of radius at most $2$ are affine, but it contains $(\bigoplus_mS_m)^2$ (Theorem~\ref{lin:movers}). Is it generated by the affine automorphisms and the maps $\Phi_\psi$, $\Phi'_\chi$ of that theorem?
\item For the affine classes of Theorem~\ref{thm:affine}, is every block map commuting with $F_f$ affine? Every automorphism is (Lemma~\ref{aff:lem:R}), and every commuting block map of radius at most $3$ is (Remark~\ref{aff:rem:endo}).
\item In every example of this paper with a nonlinear rule, the automorphisms outside $\langle\sh,F\rangle$ come from a blocking word with a degree of freedom next to the blocking column. For a rule that is not affine, the absence of blocking words does not force the automorphism group to be virtually $\Z^2$: the class of $95$ is a counterexample \cite{PaperE}. Question~\ref{q:dichotomy} asks whether it is the only one among the elementary rules.
\end{enumerate}
